\section{Introduction}
\label{sec:intro}

Fractional-order controllers promise finer tuning than their integer-order
counterparts, at the price of extra degrees of freedom that are readily handed
to metaheuristics. For a power converter, this promise meets three
implementation realities: sampling and delay, actuator saturation and
anti-windup, and the rational approximation of fractional operators.

This study starts from a documented failure. An FO-LQI $\to$ FO-PIDF campaign
on a non-ideal 48~V Buck converter spent more than 1800 evaluations across six
algorithms without ever reaching robust feasibility. The diagnosis retained at
the time blamed clipping of the projected proportional gain at the bound
$\log_{10}k_p=-4$, reached by $79~\%$ of the candidates.

\paragraph{Methodological question.}
When a parametric search fails, how can one tell whether the failure comes
from the algorithm, from the parameterisation, or from the problem itself? We
answer with a three-step protocol: a bound-saturation audit that separates
structural from pathological saturation; joint probes in controller space,
which test whether the feasible region is reachable independently of the
parameterisation; and a preregistered confirmatory campaign that measures the
failure with appropriate statistics instead of assuming it.

\paragraph{Contributions.}
\begin{enumerate}
\item A bound-saturation audit of the parameterisation, which shows here that
  clipping of $k_p$ is not the cause of the earlier failure.
\item A reconstructed and verified FO-LQI $\to$ FO-PIDF chain: the projection
  reproduces the exact integer-order identity and measures its error outside
  that domain.
\item A preregistered three-arm confirmatory campaign whose negative results
  are reported in full, sterile seeds included.
\item The identification of a controller-independent physical obstruction
  (the voltage floor at minimum duty ratio) and of a blind spot in the
  specification (limit cycles invisible to the constraints).
\end{enumerate}
